\section{Barrier geometry and the comparison problem}
\label{sec:foundations}

The paper makes three different comparisons: a central arc with the
shortest path to the same endpoint; a central arc with the shortest path
to any point that is at least as accurate; and primal movement with the
movement of a feasible primal--dual pair. Each statement specifies its
set of endpoints and its metric.

\subsection{Definitions and step counts}
Let $\Omega$ be an open convex subset of a finite-dimensional real inner
product space $E$. A $C^3$ function $F:\Omega\to\R$ with positive definite
Hessian defines a local norm, curve length, and distance:
\begin{equation}
 \norm{h}_{F,x}^2=D^2F(x)[h,h],\qquad
 \Len_F(\gamma)=\int_0^1\norm{\dot\gamma(t)}_{F,\gamma(t)}\dd t,
 \qquad
 d_F(x,y)=\inf_{\gamma:x\leadsto y}\Len_F(\gamma).
 \label{eq:metric-definitions}
\end{equation}
The infimum is over absolutely continuous curves in $\Omega$.
On an affine slice of $\Omega$, norms are restricted to the tangent
space of the slice, and curves stay in the slice. Complex matrix
spaces are treated as real vector spaces, with inner product $\operatorname{Re}\tr(X^*Y)$.

We call $F$ a standard self-concordant barrier if it tends to $+\infty$
at every finite boundary point and satisfies
\begin{equation}
 \abs{D^3F(x)[h,h,h]}\leq2\norm{h}_{F,x}^3.
 \label{eq:self-concordance}
\end{equation}
Its barrier parameter is any $\nu$ satisfying
$\abs{DF(x)[h]}\leq\sqrt\nu\norm{h}_{F,x}$ for all $x,h$; the least
such $\nu$ is the exact parameter.
In ``standard self-concordant'', the word ``standard'' refers to the
constant $2$ in \eqref{eq:self-concordance}; ``the standard barriers''
means the logarithmic and log-determinant barriers of
\cref{sec:exact-distance} and their products, with any positive scales.
They are standard self-concordant when every scale is at least one
(\cref{prop:standard-parameters}).
Self-concordance implies the polarized inequality
\begin{equation}
 \abs{D^3F(x)[u,h,h]}\leq2\norm{u}_{F,x}\norm{h}_{F,x}^2.
 \label{eq:mixed-SC}
\end{equation}
We write $\norm{h}_x$ when $F$ is fixed.

\begin{lemma}[Distance and the local norm of a chord]
\label{lem:chord-distance}
For a standard self-concordant barrier,
\begin{equation}
 \log(1+\norm{y-x}_x)\leq d_F(x,y).
 \label{eq:global-chord}
\end{equation}
If $q=\norm{y-x}_x<1$, the straight segment from $x$ to $y$ has length
at most $-\log(1-q)$. In particular, for $0<R<1$,
\begin{equation}
 \norm{y-x}_x\leq R\ \Longrightarrow\ d_F(x,y)\leq-\log(1-R),
 \qquad
 d_F(x,y)\leq\log(1+R)\ \Longrightarrow\ \norm{y-x}_x\leq R.
 \label{eq:two-chord-bounds}
\end{equation}
\end{lemma}
\begin{proof}
Along any curve $\gamma$ starting at $x$, let
$\ell(t)=\int_0^t\norm{\dot\gamma(u)}_{\gamma(u)}\dd u$.
For a fixed vector $h\ne0$, \eqref{eq:mixed-SC} gives
$\abs{\frac{\mathrm d}{\mathrm dt}\log\norm{h}_{\gamma(t)}}
\leq\dot\ell(t)$ almost everywhere. Thus
$\norm{h}_x\leq e^{\ell(t)}\norm{h}_{\gamma(t)}$ for every $h$.
The triangle inequality in the fixed norm at $x$ now gives
\[
 \norm{y-x}_x\leq\int_0^1\norm{\dot\gamma(t)}_x\dd t
 \leq\int_0^1e^{\ell(t)}\dot\ell(t)\dd t=e^{\Len_F(\gamma)}-1.
\]
Taking the infimum proves \eqref{eq:global-chord}.
For the segment $x+th$, $h=y-x$, set $a(t)=\norm{h}_{x+th}$.
Self-concordance gives $\abs{a'(t)}\leq a(t)^2$ and hence
$a(t)\leq q/(1-tq)$ while $tq<1$. Integration from $0$ to $1$ gives
the second assertion. The segment lies in $\Omega$ by convexity.
\end{proof}

These estimates underlie the short-step counts of Nesterov and
Todd~\cite[Section~3]{NesterovTodd2002}.
For a target set $A\subset\Omega$ and a starting point $x_0$, let
$D_F(x_0,A)=\inf_{x\in A}d_F(x_0,x)$.
A forward-$R$ sequence is a finite sequence $x_0,x_1,\ldots,x_T$ of
points in $\Omega$ such that $\norm{x_{j+1}-x_j}_{x_j}\leq R$; its $T$
steps are called forward $R$-chords.

\begin{corollary}[Optimal bounded movement]
\label{cor:optimal-movement}
Let $F$ be a standard self-concordant barrier and $0<R<1$.
Every forward-$R$ sequence from $x_0$ to $A$ satisfies
\begin{equation}
 T\geq\frac{D_F(x_0,A)}{-\log(1-R)}.
 \label{eq:movement-lower}
\end{equation}
If a shortest curve to $A$ exists, there is such a sequence with
\begin{equation}
 T\leq\left\lceil\frac{D_F(x_0,A)}{\log(1+R)}\right\rceil.
 \label{eq:movement-upper}
\end{equation}
\end{corollary}
\begin{proof}
For the lower bound, join consecutive points by straight segments and
apply \cref{lem:chord-distance}. For the upper bound, divide a shortest
curve into arcs of length at most $\log(1+R)$ and apply
\eqref{eq:global-chord} to their endpoints. If the distance is zero and
attained, no step is needed.
\end{proof}

For a bounded domain and a nonzero objective vector $c\in E$, define the
support value, the objective gap, and the $\eps$-accurate set:
\begin{equation}
 h_\Omega(c)=\sup_{x\in\Omega}\ip{c}{x},\qquad
 g_c(x)=h_\Omega(c)-\ip{c}{x},\qquad
 A_\eps=\{x\in\Omega:g_c(x)\leq\eps\}.
 \label{eq:gap-target}
\end{equation}
When $F$ is a strictly convex barrier, the central point with parameter
$\eta$ is
\begin{equation}
 x(\eta)=\argmin_{x\in\Omega}\{F(x)-\eta\ip{c}{x}\},\qquad \eta\geq0.
 \label{eq:central-path}
\end{equation}
At $\eta=0$ it is the analytic center $x_0$.
We write $L_\opt(\eps)=D_F(x_0,A_\eps)$ and
$L_\CP(\eta)=\Len_F(x|_{[0,\eta]})$; $L_\CP(\eta_0,\eta_1)$ is the
length of the arc between parameters $\eta_0$ and $\eta_1$, and
$L_\CP(\eps)=L_\CP(\eta_\eps)$, where $\eta_\eps$ is the first
parameter with $g_c(x(\eta_\eps))=\eps$. We say whether the argument
of $L_\CP$ is a parameter or an accuracy whenever this is ambiguous.
